\begin{table}[htbp]
\centering
\caption{Profile cụm FCM thường và FCM có trọng số do PSO tìm trên Việt Nam (c = 3, seed 1000): trung bình giá trị gốc}
\label{tab:e4-profile-viet_nam}
\footnotesize
\setlength{\tabcolsep}{3pt}
\begin{tabular}{@{}lrrrrrrrrr@{}}
\toprule
Phương pháp & Cụm & Cỡ (\%) & Vỡ nợ (\%) & Mơ hồ (\%) & \rotatebox{60}{X17} & \rotatebox{60}{X11} & \rotatebox{60}{X6} & \rotatebox{60}{X10} & \rotatebox{60}{X8} \\
\midrule
FCM (k-means++) &1 & 40,9 & 1,3 & 100,0 & 2,09 & 0,19 & 0,24 & 0,26 & 0,22 \\
 &2 & 24,8 & 2,5 & 100,0 & 0,47 & 0,36 & 0,24 & 0,58 & 0,54 \\
 &3 & 34,3 & 35,4 & 90,7 & 0,39 & 0,67 & 1,28 & 2,54 & 2,13 \\
\midrule
PSO trọng số &1 & 56,6 & 3,1 & 17,6 & 0,51 & 0,34 & 0,29 & 0,65 & 0,60 \\
 &2 & 23,7 & 7,6 & 20,4 & 3,26 & 0,23 & 0,29 & 0,42 & 0,32 \\
 &3 & 19,6 & 49,8 & 26,7 & 0,24 & 0,73 & 1,86 & 3,34 & 2,74 \\
\bottomrule
\end{tabular}
\par\smallskip\parbox{\linewidth}{\scriptsize Ghi chú: cụm đánh số theo tỉ lệ vỡ nợ tăng dần; mơ hồ: tỉ lệ điểm trong cụm có độ thuộc lớn nhất $<$ 0,6; thuộc tính chọn là 5 thuộc tính có khoảng trung bình giữa các cụm (PSO trọng số) rộng nhất tính theo độ lệch chuẩn. X17: P/S; X11: phải trả người bán/tổng tài sản; X6: tổng nợ/VCSH; X10: phải trả người bán/VCSH; X8: nợ ngắn hạn/VCSH.}
\end{table}
